\section{A Taxonomy of Autoresearch Agent Failure Patterns}
\label{sec:taxonomy}


Applying this evaluation across the collected trajectories yields \taxonomy, which organizes the 45 observed failure patterns along two orthogonal axes. The \textbf{stage axis} locates \emph{where} a failure manifests in the research pipeline---ideation, retrieval and synthesis, execution, analysis, writing, and review (Stages A--F). The \textbf{root-cause axis} identifies the underlying mechanism---\emph{why} it occurs rather than where---so that superficially different failures sharing a cause are grouped together, while one stage's failures can split across causes. A cross-stage layer (\textbf{X}) captures dynamic failures that do not localize to any single stage but describe how errors propagate, drift, or compound across the pipeline (e.g., error propagation, goal drift). A trajectory may carry multiple (stage, root-cause) instances; the judge emits one label per detected instance. This section presents the taxonomy's structure and the payoff of the two-axis view; the full label set, and definitions are in Appendix~\ref{app:failure_pattern_definitions}.



\begin{table*}[htbp]
\centering
\small
 \caption{\textbf{Systemic root-cause classification of AutoResearch failure patterns.}
  All 45 patterns are grouped under four root-cause pillars---Grounding \& Faithfulness,
  Cognitive Depth \& Adaptability, \add{Scientific} Integrity \& Alignment, and Engineering Robustness---that
  capture the underlying mechanism of failure rather than the pipeline stage at which it surfaces.
  Each pattern ID encodes its stage (with the X series spanning multiple stages), and every pattern
  maps to exactly one pillar.}
  \label{tab:root_cause_mapping}
\begin{tabularx}{\textwidth}{l X X}
\toprule
\textbf{Root Cause Pillar} & \textbf{Core Failure Focus} & \textbf{Mapped Failure Patterns (IDs)} \\
\midrule
\textbf{\add{R}1. Grounding \& Faithfulness}
& Disconnect between high-level claims/hypotheses and ground-truth code, data, or logs.
& A.6, B.1, B.2, B.5, C.3, D.1, D.4, D.6, E.1, E.4, F.6, X.6 \\
\midrule
\textbf{\add{R}2. Cognitive Depth \& Adaptability}
& Shallow reasoning/search, passivity in self-review, and inability to re-plan or pivot.
& A.1, A.3, B.4, B.6, C.6, C.7, D.5, F.1, F.2, F.3, F.4, X.3, X.7 \\
\midrule
\textbf{\add{R}3. Integrity \& Alignment}
& Metric hacking, shortcut reliance, confirmation bias, overclaiming, and goal drift.
& A.2, A.5, C.1, C.2, D.2, D.3, D.7, E.2, E.3, F.5, X.2, X.4, X.5 \\
\midrule
\textbf{\add{R}4. Engineering Robustness}
& Numerical overflows, unhandled runtime errors, and broken interaction with CLI/OS.
& A.4, B.3, C.4, C.5, C.8, X.1, X.8 \\
\bottomrule
\end{tabularx}

\end{table*}
 
\paragraph{Stage axis (A--F) and cross-stage layer (X).}
The six stages and the cross-cutting layer group the failure patterns as follows:
\begin{itemize}[leftmargin=1.6em, itemsep=2pt, parsep=0pt]
    \item \textbf{A $\cdot$ Ideation \& Planning} (6 patterns): failures of hypothesis formation and experimental design, e.g., frame-lock in a narrow hypothesis space (A.1), unfalsifiable hypotheses (A.2), and experiments that do not actually test the stated hypothesis (A.6).
    \item \textbf{B $\cdot$ Retrieval \& Synthesis} (6 patterns): failures of evidence acquisition and use, e.g., hallucinated citations (B.1), shallow search coverage (B.4), and retrieved knowledge that never informs experimental design (B.2).
    \item \textbf{C $\cdot$ Execution \& Implementation} (8 patterns): failures during coding and experimentation, e.g., circular validation (C.1), grader-fitting and data leakage (C.2), and code that diverges from the claimed methodology (C.3).
    \item \textbf{D $\cdot$ Analysis \& Interpretation} (7 patterns): failures of inference from results, e.g., mistaking artifacts for insights (D.1), confirmation bias (D.2), and fabricated metrics or tables (D.6).
    \item \textbf{E $\cdot$ Writing \& Documentation} (4 patterns): failures of faithful reporting, e.g., claims untraceable to actual execution (E.1) and overclaiming with concealed negative results (E.2).
    \item \textbf{F $\cdot$ Self-Verification \& Review} (6 patterns): failures of the agent's own quality gate, e.g., superficial checklist-style self-review (F.1) and review score hacking (F.5).
    \item \textbf{X $\cdot$ Cross-Stage Patterns} (8 patterns): dynamic failures spanning stages, e.g., cascading error propagation (X.1), goal drift (X.2), and right-for-the-wrong-reason successes (X.6).
\end{itemize}

\begin{tcolorbox}[colback=blue!5!white,colframe=blue!75!black,
  title=\textbf{Ultimate Root Cause: Metacognitive Loop}]
While failure patterns span all stages of the research lifecycle, they
are fundamentally driven by a single overarching limitation:
\textbf{Metacognitive Loop}. A human researcher works in a closed
loop---recognising the limits of what they know, checking whether
intermediate results hold up, and  re-planining when it is necessary.
Current autoresearch agents lack this loop. They can execute each
step of the research process, but they do not have the awareness to monitor what they have
produced, judge whether it is valid, and re-plan when it is
not.
\end{tcolorbox}
 
\paragraph{Root Cause axis.}
The second axis exposes structure invisible to a stage-only list: when patterns are aligned by root mechanism, families emerge that pipeline position would otherwise scatter. This core limitation manifests through four practical root-cause pillars: \textbf{Grounding \& Faithfulness}, \textbf{Cognitive Depth \& Adaptability}, \textbf{Scientific Integrity \& Alignment}, \textbf{Engineering Robustness}. We will introduce more in \S\ref{sec:rootcause}.
 
Table~\ref{tab:root_cause_mapping} gives the full pattern-to-pillar mapping. The most striking family the table reveals is the Depth root-cause: a single cognitive deficit resurfaces at nearly every stage of the pipeline, from ideation-time frame-lock (A.1) through execution-time local optimization (C.6) to review-time passivity (F.1--F.4).

 
%\vspace{0.8em}
 
\medskip
 

 
We defer empirical evidence and detailed case studies to Section~\ref{sec:analysis} and Appendix~\ref{app:detailed case studies}.



